\section{Corruption Experiment：模型是否真的使用其他 datapoints}

Benchmark performance 不能直接证明跨 datapoint attention 发挥了作用：一个高容量 NPT 可能学会忽略其他 rows，退化成普通 per-row parametric model。作者因此设计 intervention-like test，在保持当前 target row 不变的同时，破坏所有其他 rows 的关系信息，再观察预测是否恶化。

\lecturefigure{slide-24-corruption-goal.jpg}{Corruption experiment 的目标：禁止模型从其他 datapoints 获得有用关系。}{Stanford Online 官方视频 00:51:16--00:51:31。}

对每个被评估的 target entry $(i,d)$，方法保留第 $i$ 行，并对其他 rows 的每个 attribute column 独立随机 permutation。可写成
\[
\widetilde X_{k,j}=X_{\pi_j(k),j},
\qquad k\ne i,
\]
其中每个 column $j$ 使用独立 permutation $\pi_j$。

\begin{itemize}
  \item 当前 query row $i$ 不变；
  \item 单个 column 的 marginal values 大致保留；
  \item 同一 row 内原有 feature--feature 与 feature--target correspondence 被打散；
  \item 若预测依赖其他 rows 的结构，corruption 后 performance 应下降。
\end{itemize}

\lecturefigure{slide-25-corruption-method.jpg}{Corruption method：逐 attribute 打乱其他 rows，保留 query row。}{Stanford Online 官方视频 00:51:32--00:53:01。}

\teachervoice{讲者特别说明不想用简单噪声把输入“弄坏”，因为模型可能只检测到 distribution shift。逐 column permutation 尽量保留 marginal statistics，却摧毁跨字段、跨样本的可用对应关系。}

\subsection{结果：有些任务强依赖，有些任务主动忽略}

若 NPT 在所有数据集都同样依赖其他 rows，corruption 应普遍造成巨大损失；实际结果更细致：Protein 等任务显著恶化，Forest、Kick、Breast Cancer 等几乎不变。这说明 architecture 提供了 relation channel，但 optimization 可以按数据决定是否使用。

\lecturefigure{slide-26-corruption-results.jpg}{破坏其他 datapoints 信息后的 performance change。}{Stanford Online 官方视频 00:53:02--00:54:43。}

\begin{table}[H]
\centering
\small
\caption{论文 Table 2：corruption 后的 performance 变化；负值代表变差。}
\begin{tabularx}{\textwidth}{L{0.21\textwidth}*{4}{>{\centering\arraybackslash}X}}
\toprule
分类数据集 & CIFAR-10 & Poker & Income & Higgs \\
\midrule
$\Delta$ Accuracy & $-1.2$ & $-1.1$ & $-1.1$ & $-0.5$ \\
\bottomrule
\end{tabularx}
\vspace{0.6em}
\begin{tabularx}{\textwidth}{L{0.21\textwidth}*{4}{>{\centering\arraybackslash}X}}
\toprule
分类数据集 & MNIST & Forest & Kick & Breast Cancer \\
\midrule
$\Delta$ Accuracy & $-0.4$ & $-0.1$ & $-0.1$ & $0.0$ \\
\bottomrule
\end{tabularx}
\vspace{0.6em}
\begin{tabularx}{\textwidth}{L{0.21\textwidth}*{4}{>{\centering\arraybackslash}X}}
\toprule
回归数据集 & Yacht & Protein & Boston & Concrete \\
\midrule
$\Delta\mathrm{RMSE}/\mathrm{RMSE}$ & $-52\%$ & $-21\%$ & $-20\%$ & $-7\%$ \\
\bottomrule
\end{tabularx}
\end{table}

\begin{importantbox}{更有意思的不是“全部下降”}
在某些 datasets 上，NPT 发现其他 rows 没有足够价值，于是可以学会忽略它们并退化为近似 parametric predictor；在另一些 datasets 上，corruption 使 performance 严重崩溃。可学习的关键不是固定采用 non-parametric rule，而是选择依赖程度。
\end{importantbox}

\begin{warningbox}{Corruption 支持机制依赖，但不是完整因果证明}
Performance drop 表明预测利用了被破坏的跨-row structure；它不能单独告诉我们某个 attention head 的唯一因果路径，也不能推出“non-parametric 普遍优于 parametric”。Corruption 同时改变了高阶联合分布，仍可能引入模型敏感的 out-of-distribution pattern。
\end{warningbox}

\teachervoice{论文与讲座都谨慎指出：Forest、Kick、Breast Cancer 几乎不受影响，并非实验失败，而是模型可能判断跨 datapoint lookup 没有优势。这种可退化性恰好支持“依赖程度由数据学习”的主张。}

\subsection{本章小结}

Corruption test 从“模型有跨行 attention”推进到“模型在若干真实任务中实际依赖跨行信息”。证据来自反事实式的 performance change，但实验不能识别唯一机制，也不能泛化成 parametric 与 non-parametric 的普遍优劣结论。

